\documentclass{article}

% Language setting
% Replace `english' with e.g. `spanish' to change the document language
\usepackage[english]{babel}

% Set page size and margins
% Replace `letterpaper' with `a4paper' for UK/EU standard size
\usepackage[letterpaper,top=2cm,bottom=2cm,left=3cm,right=3cm,marginparwidth=1.75cm]{geometry}

% Useful packages
\usepackage[utf8]{inputenc}
\usepackage[T1]{fontenc}
\usepackage{amsmath}
\usepackage{amssymb}
\usepackage{graphicx}
\usepackage{booktabs}
\usepackage[colorlinks=true, allcolors=blue]{hyperref}
\usepackage{cite}
\usepackage{algorithm}
\usepackage{algpseudocode}
\usepackage{caption}
\usepackage{subcaption}
\usepackage{multirow}
\usepackage{array}
\usepackage{enumitem}

\title{
  \Large \textbf{Rugby Video Analytics: \\
  AI-Based Tackle Analysis Robust to Keypoint Occlusion} \\
  \large Capstone Project --- Interim Report
}
\author{
  Yu Araki \\
  \small A0304700X \\
  National University of Singapore
}
\date{\today}



\begin{document}
\maketitle
\tableofcontents
\newpage

\section{Introduction}

In modern sports, to keep winning in a competitive situation, every sport requires not only physical skill but also data analytics. Video analytics, 
particularly through human pose estimation, has emerged as a powerful 
tool in sports science, enabling automated analysis of athletes' 
movements to optimize performance and prevent injury~\cite{zheng2021survey}. 
Conventional coaching relies on subjective visual assessment by human 
experts, which is limited by the observer's experience, attention span, 
and inability to quantify subtle kinematic differences. Recent advances 
in deep learning-based pose estimation have opened the possibility of 
analyzing standard video footage by extracting joint 
coordinates~\cite{zheng2021survey}.

Rugby football, as a collision-driven sport, presents significant 
challenges for video-based motion analytics. Unlike most other sports, 
rugby involves frequent physical collisions such as tackles where appropriate body posture is not just a performance 
factor but a safety-critical requirement. Precise kinematic analysis 
of collision posture is therefore essential. However, current 
state-of-the-art models, including MediaPipe~\cite{lugaresi2019mediapipe} 
and YOLO-Pose~\cite{maji2022yolopose}, exhibit significant reliability 
degradation in such scenarios, as overlapping bodies cause severe 
occlusion of body keypoints~\cite{ji2025gct}.

This project aims to develop an AI-based coaching pipeline for rugby 
tackle analysis that is robust to keypoint occlusion. In the first 
phase of this capstone project, inspired by the DETR + GCT (Graph Convolutional Transformer) architecture proposed by Ji et al.~\cite{ji2025gct}---which combines Detection Transformers with 
graph-based keypoint refinement and gating mechanisms---we propose a YOLO\,+\,GCT pipeline. In this approach, YOLO-Pose serves as the base keypoint detector, and a gating mechanism identifies low-confidence keypoints that are passed to a GCT module for spatial refinement. The 
GCT module estimates occluded or low-confidence joint positions from neighboring visible joints via graph 
convolution~\cite{kipf2017gcn,dwivedi2021graphtransformer} and 
multi-head self-attention~\cite{vaswani2017attention}.

This interim report presents the results of Phase~1 and addresses 
the limitations that motivate the Phase~2 plan: extending the spatial 
refinement model to incorporate temporal dependencies across video frames.

% =====================
\section{Problem Definition}
% =====================

Given a video sequence $\mathbf{V} = \{I_1, I_2, \ldots, I_T\}$ consisting
of $T$ frames, the task of human pose estimation is to predict, for each
frame $I_t$, a set of $K$ body keypoints.
Each keypoint $k \in \{1, 2, \ldots, K\}$ is represented as a tuple:
\begin{equation}
  \hat{p}_k^{(t)} = \bigl(\hat{x}_k^{(t)},\; \hat{y}_k^{(t)},\; c_k^{(t)}\bigr),
  \label{eq:keypoint}
\end{equation}
where $(\hat{x}_k^{(t)}, \hat{y}_k^{(t)})$ denote the predicted 2-D pixel
coordinates of keypoint $k$ at frame $t$, and $c_k^{(t)} \in [0,1]$ is the
associated confidence score reflecting the model's certainty in that
prediction.
In the COCO keypoint format adopted by YOLO-Pose~\cite{maji2022yolopose},
$K = 17$ joints are defined, covering the nose, eyes, ears, shoulders,
elbows, wrists, hips, knees, and ankles.

In rugby tackle situations, physical contact between the tackler and the
ball carrier causes partial or full occlusion of body parts.
A keypoint $k$ at frame $t$ is defined as \emph{occluded} when its
confidence score falls below a threshold $\tau$:
\begin{equation}
  \mathcal{O}(t)
    = \bigl\{\, k \in \{1,\ldots,K\} \;\big|\; c_k^{(t)} < \tau \,\bigr\},
  \label{eq:occlusion_set}
\end{equation}
where $\mathcal{O}(t)$ denotes the set of occluded keypoints at frame $t$.
In this study, $\tau = 0.5$ is adopted, consistent with the standard COCO
evaluation threshold.

The occlusion problem manifests in two failure modes that reduce the
reliability of downstream coaching analysis:

\begin{enumerate}

  \item \textbf{Positional jitter}~\cite{zeng2022smoothnet}.
  When $c_k^{(t)}$ is low, the predicted coordinates
  $(\hat{x}_k^{(t)},\hat{y}_k^{(t)})$ exhibit large frame-to-frame
  fluctuations.
  This is quantified by the \emph{jitter metric} $J_k$, defined as the
  mean Euclidean displacement between consecutive frames for keypoint $k$:
  \begin{equation}
    J_k = \frac{1}{T-1}
          \sum_{t=2}^{T}
          \left\|
            \hat{p}_k^{(t)} - \hat{p}_k^{(t-1)}
          \right\|_2,
    \label{eq:jitter}
  \end{equation}
  where $\|\cdot\|_2$ denotes the Euclidean norm.
  A high $J_k$ value during the collision phase indicates unreliable
  tracking of joint $k$, which directly corrupts any downstream
  joint-angle or angular-velocity computation.

  \item \textbf{Identity switching}.
  In multi-person scenarios, detectors can re-assign keypoints from the target athlete to a nearby
  individual when bodies overlap~\cite{zheng2021survey}.
  This is particularly problematic for tackle analysis, where the tackler
  and ball carrier are in close proximity, and the detector must
  consistently track a specific individual throughout the sequence.

\end{enumerate}

Given the above, the main research question of this project is:
\begin{quote}
  \textit{How can positional jitter $J_k$ for occluded keypoints
  ($k \in \mathcal{O}(t)$) be reduced during rugby tackle sequences,
  while maintaining the accuracy of non-occluded keypoints?}
\end{quote}

Specifically, the objective is to design a spatial refinement module that
takes the raw YOLO-Pose output
$\bigl\{\hat{p}_k^{(t)}\bigr\}_{k=1}^{K}$
as input and produces a refined estimate
$\bigl\{\tilde{p}_k^{(t)}\bigr\}_{k=1}^{K}$ such that:
\begin{equation}
  J_k^{\,\text{refined}} < J_k^{\,\text{raw}},
  \qquad \forall\, k \in \mathcal{O}(t).
  \label{eq:objective}
\end{equation}
The design of the refinement module and its confidence-aware blending
strategy are detailed in Section~\ref{sec:methodology}.

\section{Literature Review}
\label{sec:literature}

This section reviews the key technical components that underpin the proposed
pipeline: per-frame pose estimation models (Section~\ref{sec:lit:pose}),
graph-based skeleton modelling with GCN (Section~\ref{sec:lit:gcn}),
and the Graph Convolutional Transformer with gating mechanisms
(Section~\ref{sec:lit:gct}).

% ---------------------------------------------------------
\subsection{Per-Frame Pose Estimation: MediaPipe and YOLO-Pose}
\label{sec:lit:pose}
% ---------------------------------------------------------

Both MediaPipe~\cite{lugaresi2019mediapipe} and YOLO-Pose~\cite{maji2022yolopose}
operate on a per-frame basis: each video frame is processed independently to
predict the 2-D coordinates $(\hat{x}_k^{(t)}, \hat{y}_k^{(t)})$ and an
associated confidence score $c_k^{(t)}$ for each of $K$ body keypoints.

\paragraph{MediaPipe.}
MediaPipe Pose~\cite{lugaresi2019mediapipe} is a lightweight, real-time
pose estimation framework developed by Google.
It operates in a two-stage pipeline: a person detector first localises
the bounding box of the subject, and a landmark regression network
(BlazePose~\cite{bazarevsky2020blazepose}) then predicts 33 body
keypoints within the detected region.
Each keypoint is represented as $(x, y, z, v)$, where
$(x, y) \in [0,1]$ are normalised image coordinates, $z$ is a relative
depth estimate, and $v \in [0,1]$ is a visibility score indicating the
likelihood that the joint is visible in the frame.

\paragraph{YOLO-Pose.}
YOLO (You Only Look Once) is a family of single-shot object detection
models that process an entire image in one forward pass, jointly predicting
bounding boxes and class probabilities~\cite{maji2022yolopose}.
YOLO-Pose extends this architecture to simultaneously output $K = 17$ body
keypoints in the COCO format for each detected person.
For each keypoint $k$, the model outputs
$(\hat{x}_k^{(t)},\, \hat{y}_k^{(t)},\, c_k^{(t)})$,
where $(\hat{x}_k^{(t)}, \hat{y}_k^{(t)})$ are pixel coordinates and
$c_k^{(t)} \in [0,1]$ is a confidence score.

A key architectural difference from MediaPipe is that YOLO-Pose performs
person detection and keypoint regression in a single unified network, whereas
MediaPipe uses a two-stage cascade.
This makes YOLO-Pose inherently capable of handling multi-person scenarios,
as each detected bounding box carries its own independent set of keypoints---a
property that is essential for tackle analysis involving two players.

% ---------------------------------------------------------
\subsection{Graph Convolutional Networks for Skeleton Modelling}
\label{sec:lit:gcn}
% ---------------------------------------------------------

The human skeleton forms a natural graph structure: joints serve as nodes and
bones serve as edges.
This observation motivates the use of Graph Convolutional Networks (GCN)
for pose refinement~\cite{kipf2017gcn}.
Given $K = 17$ keypoints, a GCN defines an undirected graph
$\mathcal{G} = (\mathcal{V}, \mathcal{E})$,
where each vertex $v_k \in \mathcal{V}$ represents a body joint and each
edge $(v_i, v_j) \in \mathcal{E}$ encodes an anatomical connection.
For example, if a wrist keypoint is occluded, its position can be inferred
from the visible elbow via graph connectivity, as the wrist is anatomically
``downstream'' of the elbow along the kinematic chain.

The skeleton connectivity is encoded in an adjacency matrix
$A \in \mathbb{R}^{K \times K}$, where $A_{ij} = 1$ if joints $i$ and $j$
are connected and $A_{ij} = 0$ otherwise.
To allow each node to retain its own features during aggregation, self-loops
are added, giving $\hat{A} = A + I$.
To prevent feature magnitudes from scaling with node degree---a problem that
arises because high-degree nodes such as the shoulder aggregate information
from more neighbours than low-degree nodes such as the wrist---the adjacency
matrix is symmetrically normalised following Kipf \&
Welling~\cite{kipf2017gcn}:
\begin{equation}
  \tilde{A} = \hat{D}^{-1/2}\,\hat{A}\,\hat{D}^{-1/2},
  \label{eq:norm_adj}
\end{equation}
where $\hat{D}$ is the diagonal degree matrix with
$\hat{D}_{ii} = \sum_{j} \hat{A}_{ij}$.
Each element of $\tilde{A}$ is scaled by $1/\!\sqrt{d_i d_j}$, which
ensures that both high-degree and low-degree nodes are updated at the
same magnitude.

The GCN layer update rule for node features
$H^{(\ell)} \in \mathbb{R}^{K \times d_\ell}$ at layer $\ell$ is then:
\begin{equation}
  H^{(\ell+1)} = \sigma\!\left(\tilde{A}\,H^{(\ell)}\,W^{(\ell)}\right),
  \label{eq:gcn_update}
\end{equation}
where $W^{(\ell)} \in \mathbb{R}^{d_\ell \times d_{\ell+1}}$ is a
learnable weight matrix and $\sigma(\cdot)$ is a nonlinear activation
function (ReLU in this work).
Equation~\eqref{eq:gcn_update} can equivalently be written node-wise as
\begin{equation}
  h_i^{(\ell+1)} = \sigma\!\left(
    \sum_{j \in \mathcal{N}(i)} \frac{1}{\sqrt{d_i\, d_j}}\,
    W^{(\ell)}\, h_j^{(\ell)}
  \right),
  \label{eq:gcn_nodewise}
\end{equation}
which shows explicitly how each joint aggregates information from its
anatomical neighbours $\mathcal{N}(i)$, weighted by the inverse square
root of the degree product.

% ---------------------------------------------------------
\subsection{Graph Convolutional Transformer with Gating Mechanisms}
\label{sec:lit:gct}
% ---------------------------------------------------------

While GCN is effective at propagating information between anatomically adjacent
joints, a single GCN layer aggregates features only from direct neighbours.
To capture global dependencies---such as the relationship between a wrist and
the opposite hip during a tackle---Dwivedi and Bresson~\cite{dwivedi2021graphtransformer}
proposed the Graph Transformer (GT), which augments graph convolution with
multi-head self-attention (MHA)~\cite{vaswani2017attention}.
Ji et al.~\cite{ji2025gct} subsequently integrated this architecture into a
pose estimation pipeline as a Graph Convolutional Transformer (GCT), combined
with a gating mechanism for confidence-aware feature filtering.

\paragraph{Multi-Head Self-Attention.}
MHA computes attention over all $K$ nodes simultaneously.
Given node features $H \in \mathbb{R}^{K \times d}$, the query, key, and
value matrices are projected as
$Q = H W_Q$, $K_{\text{mat}} = H W_K$, $V = H W_V$, and the
attention output is:
\begin{equation}
  \text{Attention}(Q, \mathbf{K}, V) = \text{softmax}\left(\frac{Q\mathbf{K}^\top}{\sqrt{d_k}}\right)V
  \label{eq:mha}
\end{equation}
where $d_k$ is the dimension of each attention head and the factor
$1/\!\sqrt{d_k}$ stabilises the gradient by preventing the dot products from
growing large~\cite{vaswani2017attention}.
This allows every joint to directly attend to every other joint in a single
layer, overcoming the locality constraint of GCN.

\paragraph{Gating Mechanism.}
A gating mechanism is applied to selectively suppress features of
low-confidence keypoints before residual addition~\cite{ji2025gct}.
Given an input feature vector $\mathbf{x} \in \mathbb{R}^d$, the gate
computes an element-wise scaling factor via sigmoid activation:
\begin{equation}
  \mathbf{g} = \sigma\!\left(W_g\,\mathbf{x} + \mathbf{b}_g\right),
  \label{eq:gate}
\end{equation}
where $W_g \in \mathbb{R}^{d \times d}$ is a learnable weight matrix,
$\mathbf{b}_g \in \mathbb{R}^d$ is a bias, and $\sigma(\cdot)$ is the
sigmoid function.
The gated output is:
\begin{equation}
  \mathbf{z} = \mathbf{x} \odot \mathbf{g},
  \label{eq:gated_output}
\end{equation}
where $\odot$ denotes the element-wise (Hadamard) product.
When the confidence of a keypoint is low, the gate learns to output values
near zero, suppressing unreliable features while allowing high-confidence
features to pass through unmodified.

\paragraph{GCT Layer.}
A complete GCT layer processes the input node features $H$ through three
sequential sub-layers, each with residual connections and layer
normalisation:
\begin{align}
  H'    &= \mathrm{LayerNorm}\!\left(H   + \mathrm{GCN}(H,\,\tilde{A})\right),
           \label{eq:gct1} \\
  H''   &= \mathrm{LayerNorm}\!\left(H'  + \mathrm{Gate}\!\left(
             \mathrm{MHA}(H')\right)\right),
           \label{eq:gct2} \\
  H^{\mathrm{out}}
        &= \mathrm{LayerNorm}\!\left(H'' + \mathrm{FFN}(H'')\right),
           \label{eq:gct3}
\end{align}
where $\mathrm{FFN}$ is a two-layer feed-forward network with ReLU
activation:
\begin{equation}
  \mathrm{FFN}(\mathbf{x})
    = W_2\,\mathrm{ReLU}(W_1\,\mathbf{x} + \mathbf{b}_1) + \mathbf{b}_2.
  \label{eq:ffn}
\end{equation}
The residual connections and layer normalisation stabilise training and
mitigate vanishing gradients, following standard Transformer
practice~\cite{vaswani2017attention}.
In Eq.~\eqref{eq:gct1}, GCN captures \emph{local} anatomical structure;
in Eq.~\eqref{eq:gct2}, MHA followed by gating captures \emph{global}
dependencies while suppressing low-confidence features; and
Eq.~\eqref{eq:gct3} applies a non-linear projection to the combined
representation.
Figure~\ref{fig:gct_arch} illustrates the architecture of a single GCT
layer used in this work.

\begin{figure}[htbp]
  \centering
  \includegraphics[width=0.45\textwidth]{figure/GCT_layer_architecture.png}
  \caption{Architecture of a single GCT layer: GCN (local aggregation)
           $\rightarrow$ Multi-Head Attention with Gating (global context
           and confidence filtering) $\rightarrow$ Feed-Forward Network,
           each with residual connections and layer normalisation.}
  \label{fig:gct_arch}
\end{figure}

% =============================================================================
\section{Methodology}
\label{sec:methodology}
% =============================================================================

The proposed methodology consists of a two-stage pipeline for analysing
rugby tackle videos.

% ---------------------------------------------------------
\subsection{Stage~1: Base Keypoint Detector Selection}
\label{sec:meth:stage1}
% ---------------------------------------------------------

Given an input video
$\mathbf{V} = \{I_1, \ldots, I_T\}$ containing a rugby tackle scene
with occlusion, two single-frame pose estimation models---MediaPipe Pose
and YOLO-Pose---are applied to extract keypoint predictions for each
frame.
The purpose of this stage is to compare the two models and select the
more suitable base detector for the downstream refinement pipeline.
The comparison is conducted in terms of confidence score across all 17
COCO keypoints, with particular attention to the occlusion phase of the
tackle.

% ---------------------------------------------------------
\subsection{Stage~2: GCT Spatial Refinement}
\label{sec:meth:stage2}
% ---------------------------------------------------------

The raw YOLO output
$\bigl\{\hat{p}_k^{(t)}\bigr\}_{k=1}^{K}$ is fed into a trained Graph
Convolutional Transformer (GCT), which refines the keypoint positions by
exploiting the spatial relationships of the human skeleton graph.
A confidence-aware blending strategy combines the raw YOLO output and
the GCT prediction to produce the final estimate
$\bigl\{\tilde{p}_k^{(t)}\bigr\}_{k=1}^{K}$.

\paragraph{GCT Input Features.}
For each keypoint~$k$ at frame~$t$, the input feature vector to the GCT
is a 5-dimensional vector:
\begin{equation}
  \mathbf{f}_k^{(t)} = \bigl(
    \hat{x}_k^{(t)},\;
    \hat{y}_k^{(t)},\;
    c_k^{(t)},\;
    \Delta x_k^{(t)},\;
    \Delta y_k^{(t)}
  \bigr),
  \label{eq:gct_input}
\end{equation}
where $(\hat{x}_k^{(t)}, \hat{y}_k^{(t)})$ are the normalised 2-D
coordinates predicted by YOLO, $c_k^{(t)}$ is the confidence score, and
$(\Delta x_k^{(t)}, \Delta y_k^{(t)})$ represent the frame-to-frame
velocity defined as
$\Delta x_k^{(t)} = \hat{x}_k^{(t)} - \hat{x}_k^{(t-1)}$.
The inclusion of velocity provides a simple form of temporal context
within the single-frame input.

\paragraph{Model Architecture.}
The GCT model architecture is summarised in
Table~\ref{tab:gct_hyperparams}.
The input projection maps the 5-dimensional features to a
$d_{\mathrm{model}} = 64$ dimensional hidden space.
The model consists of $L = 3$ stacked GCT layers, each containing a GCN
sub-layer, a 4-head attention sub-layer with gating, and a feed-forward
network with $d_{\mathrm{ffn}} = 128$ hidden units.
The output projection maps the hidden representation back to 2
dimensions $(x, y)$ for coordinate regression.

\begin{table}[htbp]
  \centering
  \caption{GCT model hyperparameters.}
  \label{tab:gct_hyperparams}
  \begin{tabular}{@{}ll@{}}
    \toprule
    Hyperparameter & Value \\
    \midrule
    Input features $d_{\mathrm{in}}$
      & 5 $(x, y, c, \Delta x, \Delta y)$ \\
    Hidden dimension $d_{\mathrm{model}}$ & 64 \\
    Number of GCT layers $L$ & 3 \\
    Attention heads $n_{\mathrm{heads}}$ & 4 \\
    FFN hidden dimension $d_{\mathrm{ffn}}$ & 128 \\
    Dropout rate & 0.1 \\
    Output dimension & 2 $(x, y)$ \\
    Total parameters & $\sim$126\,K \\
    \bottomrule
  \end{tabular}
\end{table}


\paragraph{Self-Supervised Training via Masked Keypoint Prediction.}
Since manually annotated ground-truth labels for rugby tackle video are
unavailable, a self-supervised training strategy based on masked keypoint
prediction is employed.
Video frames where all 17 keypoints have high confidence
($c_k > \tau_{\mathrm{GT}}$, with $\tau_{\mathrm{GT}} = 0.5$) are
selected as pseudo ground-truth.
For each such frame, $M = 20$ augmented samples are generated by
randomly masking approximately 20\% of the keypoints (3--4 keypoints per
frame) and adding Gaussian noise ($\sigma = 0.02$) to the masked
coordinates.
A separate training video of a single athlete performing agility
movements under clear visibility (no occlusion) was used, ensuring the
model learned from clean skeletal configurations.

The loss function is a masked mean squared error (Masked MSE), computed
only over the keypoints that were artificially masked:
\begin{equation}
  \mathcal{L}_{\mathrm{masked}}
    = \frac{
        \sum_{k=1}^{K} (1 - m_k) \cdot
        \bigl\| \tilde{p}_k - p_k^{\,\mathrm{GT}} \bigr\|^2
      }{
        \sum_{k=1}^{K} (1 - m_k) + \epsilon
      },
  \label{eq:masked_loss}
\end{equation}
where $m_k \in \{0, 1\}$ is the mask indicator ($m_k = 1$ if
keypoint~$k$ is visible), $\tilde{p}_k$ is the GCT prediction,
$p_k^{\,\mathrm{GT}}$ is the pseudo ground-truth coordinate, and
$\epsilon = 10^{-8}$ prevents division by zero.

The model was trained for 100 epochs using the Adam optimiser with an
initial learning rate of $10^{-3}$ and weight decay of $10^{-5}$.
A cosine annealing schedule was applied to decay the learning rate.
The batch size was set to 64, and gradient clipping with a maximum norm
of 1.0 was used to stabilise training.
An 80/20 random split of the training samples was used for training and
validation, respectively.

\paragraph{Confidence-Aware Blending.}
At inference time, the GCT is applied to every frame of the evaluation
video.
Rather than replacing the YOLO predictions entirely, a
confidence-aware blending strategy produces the final keypoint estimate:
\begin{equation}
  \tilde{p}_k^{(t)}
    = \alpha_k^{(t)} \cdot \hat{p}_k^{(t)}
      + \bigl(1 - \alpha_k^{(t)}\bigr) \cdot \hat{p}_{k,\mathrm{GCT}}^{(t)},
  \label{eq:blending}
\end{equation}
where $\alpha_k^{(t)} = \min(c_k^{(t)},\, 1)$ is the blending weight
derived from the YOLO confidence score.
When confidence is high ($c_k^{(t)} \approx 1$), the original YOLO
prediction is retained.
When confidence is low ($c_k^{(t)} \to 0$), the output shifts toward
the GCT refinement.
This ensures that the GCT does not degrade predictions in well-detected
regions while providing correction in occluded regions.


% =============================================================================
\section{Experimental Results and Analysis}
\label{sec:results}
% =============================================================================

% ---------------------------------------------------------
\subsection{Stage~1: MediaPipe vs.\ YOLO-Pose Comparison}
\label{sec:res:stage1}
% ---------------------------------------------------------

Table~\ref{tab:conf_comparison} summarises the confidence score
comparison between MediaPipe and YOLO-Pose across 126 frames of a rugby
tackle video.
YOLO achieved a higher overall mean confidence score (0.857) compared
to MediaPipe (0.722).
The most significant differences were observed for the upper limbs:
YOLO scored 0.729 and 0.688 for the left and right wrists,
respectively, compared to MediaPipe's 0.312 and 0.242.
This is particularly relevant for tackle analysis, where the arms
(typically wrapping the ball carrier) are critical to identify.
Both models showed high and comparable confidence for large-area joints
such as the hips, suggesting that the challenge lies specifically in the
distal limb and head regions where occlusion most frequently occurs.

\begin{table}[htbp]
  \centering
  \caption{Confidence/visibility score comparison: YOLO-Pose vs.\
           MediaPipe across 126 frames. Bold indicates the higher mean
           confidence per joint.}
  \label{tab:conf_comparison}
  \footnotesize
  \begin{tabular}{@{}l cc cc l@{}}
    \toprule
    Joint
      & YOLO Mean & $<$0.5\%
      & MP Mean   & $<$0.5\%
      & Winner \\
    \midrule
    nose            & 0.171 & 86.5\% & \textbf{0.963} &  3.2\% & MP   \\
    left\_shoulder  & \textbf{0.983} &  0.0\% & 0.968 &  3.2\% & Draw \\
    right\_shoulder & \textbf{0.977} &  0.0\% & 0.968 &  3.2\% & Draw \\
    left\_elbow     & \textbf{0.856} & 14.3\% & 0.485 & 54.0\% & YOLO \\
    right\_elbow    & \textbf{0.847} & 10.3\% & 0.452 & 46.0\% & YOLO \\
    left\_wrist     & \textbf{0.729} & 22.2\% & 0.312 & 76.2\% & YOLO \\
    right\_wrist    & \textbf{0.688} & 23.8\% & 0.242 & 84.9\% & YOLO \\
    left\_hip       & \textbf{0.994} &  0.0\% & 0.968 &  3.2\% & Draw \\
    right\_hip      & \textbf{0.993} &  0.0\% & 0.968 &  3.2\% & Draw \\
    left\_knee      & \textbf{0.987} &  0.0\% & 0.752 &  7.9\% & YOLO \\
    right\_knee     & \textbf{0.992} &  0.0\% & 0.712 & 10.3\% & YOLO \\
    left\_ankle     & \textbf{0.946} &  4.8\% & 0.810 &  7.9\% & YOLO \\
    right\_ankle    & \textbf{0.973} &  0.0\% & 0.787 &  5.6\% & YOLO \\
    \midrule
    \textbf{Overall}& \textbf{0.857} & 12.5\% & 0.722 & 23.7\% & \textbf{YOLO} \\
    \bottomrule
  \end{tabular}
\end{table}

Figure~\ref{fig:conf_timeline} shows the confidence scores over time for
selected joints.
Both models perform reliably during non-contact phases
($t \leq 3$\,s), with $|\mathcal{O}(t)| \approx 0$ for most frames
when $\tau = 0.5$.
During the collision phase ($t \approx 4$\,s), $|\mathcal{O}(t)|$
increases sharply, particularly for distal joints such as wrists and
elbows, with confidence scores dropping to $c_k^{(t)} < 0.3$.
The degradation is most pronounced for MediaPipe, where wrist confidence
frequently drops to zero, while YOLO maintains greater consistency
despite also experiencing reduced scores around the occluded scene.

\begin{figure}[htbp]
  \centering
  \includegraphics[width=\textwidth]{figure/confidence_timeline_plot.png}
  \caption{Confidence score over time for selected joints (right
           shoulder, hip, knee, ankle). The dashed yellow line indicates
           $\tau = 0.5$. MediaPipe shows severe instability during the
           collision phase ($t \approx 3.5$--$4.5$\,s).}
  \label{fig:conf_timeline}
\end{figure}

Based on these findings, YOLO-Pose (specifically, YOLOv11s-pose) was
selected as the base detector for the downstream GCT pipeline,
predicated on its superior mean confidence and robust multi-person
localisation capability.

% ---------------------------------------------------------
\subsection{Stage~2: GCT Refinement Results}
\label{sec:res:stage2}
% ---------------------------------------------------------

The trained GCT was evaluated on a separate rugby tackle video not used
during training.
The training video captured athletic movement under clear visibility,
while the evaluation video featured a tackle with real occlusion during
the collision phase.
Since the GCT itself does not modify the confidence score, the mean
jitter~$J_k$ (Eq.~\ref{eq:jitter}) was used as the primary evaluation
metric.
Table~\ref{tab:jitter} presents the comparison between raw YOLO and
YOLO\,+\,GCT.

\begin{table}[htbp]
  \centering
  \caption{Mean jitter comparison: YOLO vs.\ YOLO\,+\,GCT. ``Occ''
           columns are computed only over frames where
           $c_k^{(t)} < 0.5$. Positive reduction indicates improvement.}
  \label{tab:jitter}
  \footnotesize
  \begin{tabular}{@{}l cc r cc r@{}}
    \toprule
    Keypoint
      & YOLO   & GCT    & Reduce
      & Occ YOLO & Occ GCT & Occ Red \\
    \midrule
    left\_wrist     & 0.0089 & 0.0085 & +5.3\%
                    & 0.0063 & 0.0062 & +0.6\% \\
    right\_wrist    & 0.0108 & 0.0102 & +5.1\%
                    & 0.0039 & 0.0059 & $-$53.2\% \\
    left\_elbow     & 0.0063 & 0.0067 & $-$5.0\%
                    & 0.0042 & 0.0048 & $-$16.8\% \\
    right\_elbow    & 0.0072 & 0.0072 & +0.5\%
                    & 0.0000 & 0.0000 & +0.0\% \\
    left\_knee      & 0.0065 & 0.0064 & +1.3\%
                    & 0.0000 & 0.0000 & +0.0\% \\
    right\_knee     & 0.0057 & 0.0058 & $-$0.8\%
                    & 0.0000 & 0.0000 & +0.0\% \\
    nose            & 0.0070 & 0.0075 & $-$7.1\%
                    & 0.0064 & 0.0068 & $-$5.2\% \\
    left\_eye       & 0.0071 & 0.0083 & $-$17.2\%
                    & 0.0071 & 0.0083 & $-$17.2\% \\
    right\_eye      & 0.0070 & 0.0082 & $-$18.4\%
                    & 0.0069 & 0.0080 & $-$15.5\% \\
    \bottomrule
  \end{tabular}
\end{table}

Several findings emerge from the jitter analysis:

\begin{enumerate}

  \item \textbf{Marginal improvement on wrists.}
  For the joints most relevant to tackle analysis---the left and right
  wrists---the GCT achieved overall jitter reductions of 5.3\% and
  5.1\%, respectively.
  However, the right wrist jitter \emph{increased} by 53.2\% during
  occlusion-specific frames, indicating that the model remains sensitive
  to severe occlusion and can become unstable when multiple adjacent
  joints are simultaneously masked.

  \item \textbf{Many joints exhibited zero occluded frames.}
  For the majority of joints (shoulders, hips, knees), no frames were
  classified as occluded ($|\mathcal{O}(t)| = 0$), indicating the
  difficulty in evaluating GCT effectiveness under the current dataset
  and the need to refine the balance between base model and GCT output.

  \item \textbf{Face keypoints showed the worst degradation.}
  The eyes and nose exhibited the largest jitter increases ($-$17\% to
  $-$18\%).
  This is attributable to the evaluation video being recorded from
  behind the athlete, rendering face keypoints permanently invisible
  regardless of contact.
  This demonstrates the model's sensitivity to camera angle.

\end{enumerate}

% ---------------------------------------------------------
\subsection{Phase Detection Consistency}
\label{sec:res:phase}
% ---------------------------------------------------------

As an additional evaluation, angular velocity-based phase detection
results between YOLO and YOLO\,+\,GCT were compared
(Figure~\ref{fig:phase_detection}).
Both pipelines identified the same impact frame (frame~79,
$t = 2.63$\,s), with peak angular velocities of 1,737\,deg/s (YOLO) and
1,690\,deg/s (GCT).
Despite a slight reduction in peak velocity under the GCT pipeline, the
overall phase structure was preserved, confirming that the GCT
refinement does not disrupt downstream temporal analysis.

\begin{figure}[htbp]
  \centering
  \includegraphics[width=\textwidth]{figure/phase_detection_comparison.png}
  \caption{Phase detection comparison: YOLO (left) vs.\ YOLO\,+\,GCT
           (right), showing right shoulder angle (top) and angular
           velocity (bottom). Both pipelines detect the same impact
           frame at $t = 2.63$\,s.}
  \label{fig:phase_detection}
\end{figure}


% =============================================================================
\section{Conclusion}
\label{sec:conclusion}
% =============================================================================

This interim report presented a two-stage pipeline for rugby tackle
video analysis: (1)~base keypoint detector selection and
(2)~spatial refinement using a Graph Convolutional Transformer (GCT).
The Stage~1 comparison demonstrated that YOLO-Pose significantly
outperforms MediaPipe in tackle scenarios, particularly for upper-limb
keypoints critical to tackle assessment.
The Stage~2 GCT refinement produced marginal improvements for wrist
jitter ($\sim$5\%) but revealed fundamental limitations of spatial-only
refinement under severe occlusion.

Three principal factors account for these limitations:
\begin{enumerate}[label=(\alph*)]
  \item \textbf{Insufficient spatial information during severe
  occlusion.}
  When both the wrist and elbow of the same arm are simultaneously
  occluded, the GCT has no reliable anchor from which to infer the
  missing positions.
  Graph propagation from distant visible joints carries only a weak
  signal to accurately predict the precise coordinates of a distal
  joint through multiple hops.

  \item \textbf{Training--evaluation domain gap.}
  The GCT was trained on agility movement data with artificially masked
  keypoints, whereas the evaluation scenario involves real occlusion
  caused by body overlap.
  The spatial configuration of visible joints may differ significantly
  between these two settings.
  Furthermore, the model proved unstable for particular body regions
  depending on the camera angle.

  \item \textbf{Lack of temporal context.}
  The current model processes each frame independently, with only a
  single-frame velocity feature $(\Delta x, \Delta y)$ as a temporal
  cue.
  Multi-frame temporal modelling is expected to provide a much stronger
  prior for occluded keypoint estimation, as the trajectory of a joint
  across several frames constrains its plausible position far more
  effectively than a single-frame snapshot.
\end{enumerate}


% =============================================================================
\section{Challenges and Knowledge Acquired}
\label{sec:challenges}
% =============================================================================

% ---------------------------------------------------------
\subsection{Technical Challenges}
\label{sec:challenges:technical}
% ---------------------------------------------------------

The most significant technical challenge encountered in Phase~1 was the
severity of the occlusion problem in rugby tackle video.
While the literature reports high accuracy for pose estimation under
general occlusion scenarios---Ji et al.~\cite{ji2025gct} achieve a mean
Average Precision (mAP) of 93.4\% on the PoseTrack
dataset---rugby tackles present an \emph{extreme} occlusion regime:
two players' bodies are pressed together at high velocity, causing
simultaneous occlusion of multiple joints on the same kinematic chain.
This differs fundamentally from the benchmark datasets used in prior
work, where occlusion typically affects isolated joints rather than
entire limb segments.

% ---------------------------------------------------------
\subsection{Knowledge and Skills Gained}
\label{sec:challenges:skills}
% ---------------------------------------------------------

Through Phase~1, the necessary technical knowledge was acquired across
multiple domains.
In computer vision, hands-on experience with two production-grade pose
estimation frameworks---MediaPipe and YOLO-Pose---was gained, including
their different output formats, coordinate systems (normalised vs.\
pixel), and multi-person handling strategies.
In deep learning, a Graph Convolutional Transformer was implemented from
scratch in PyTorch, developing a working understanding of graph
convolution with normalised adjacency matrices, scaled multi-head
attention, sigmoid gating with Hadamard products, and self-supervised
training with masked inputs.

% =============================================================================
\section{Next Semester Plan}
\label{sec:next}
% =============================================================================

The key finding of this interim report is the limitation of a
spatial-only refinement model, even with GCT enhancement, for occluded
keypoints during rugby tackles.
This motivates the primary objective for the next semester:
incorporating \textbf{temporal correlations} across consecutive video
frames.
The hypothesis is that an occluded joint at frame~$t$ can be more
accurately estimated by leveraging its known trajectory from frames
$t\!-\!1,\, t\!-\!2,\, \ldots,\, t\!-\!w$ (where $w$ is the temporal
window size), rather than relying solely on the spatial configuration of
visible joints within a single frame.

Once the model achieves reliable keypoint tracking through occlusion,
the final stage of the project envisions integrating the pose estimation
pipeline with a Large Language Model (LLM) to generate natural-language
coaching feedback.
By converting the refined kinematic data---joint angles, angular
velocities, and body posture metrics---into structured prompts, the
system aims to deliver actionable, personalised coaching advice that
bridges the gap between quantitative biomechanical analysis and
intuitive athlete instruction.


\bibliographystyle{ieeetr}
\bibliography{sample}

\end{document}
